\section{逻辑电平、晶体管与接口契约}

第二章看到真实电压连续变化；数字逻辑把允许范围映射成有限符号。接收器只在
规定采样条件下把输入判作 0 或 1。发送器要保证输出低电平不高于接收器允许的
\(V_{IL(max)}\)，输出高电平不低于 \(V_{IH(min)}\)，并留出噪声裕量：

\[
NM_L=V_{IL(max)}-V_{OL(max)},
\qquad
NM_H=V_{OH(min)}-V_{IH(min)}.
\]

两个裕量都应为正。不同电压域之间可能需要电平转换器；把
\(\SI{5}{V}\) 输出直接接到只允许 \(\SI{3.3}{V}\) 的输入可能永久损坏器件。

\topic{布尔变量与真值表}
布尔变量只有 0 和 1。逻辑非、与、或分别写作

\[
Y=\lnot A,\qquad Y=A\land B,\qquad Y=A\lor B.
\]

真值表穷举有限输入组合：

\begin{center}
\begin{osgridtabular}{cc|ccccc}
\(A\) & \(B\) & NOT \(A\) & AND & OR & XOR & NAND \\ \hline
0 & 0 & 1 & 0 & 0 & 0 & 1 \\ \hline
0 & 1 & 1 & 0 & 1 & 1 & 1 \\ \hline
1 & 0 & 0 & 0 & 1 & 1 & 1 \\ \hline
1 & 1 & 0 & 1 & 1 & 0 & 0 \\ \hline
\end{osgridtabular}
\end{center}

XOR（异或）在输入不同的时候为 1。NAND（与非）先执行 AND 再取反。只用
NAND 就能构造其他逻辑门，因此称它为功能完备门；“完备”描述表达能力，
不表示真实芯片只由一种版图单元组成。

\begin{center}
\begin{tikzpicture}[circuit logic US,font=\footnotesize\sffamily]
  \node[not gate US,draw,logic gate inputs=nn] (not) at (0,0) {};
  \node[and gate US,draw,logic gate inputs=nn] (and) at (3,0) {};
  \node[or gate US,draw,logic gate inputs=nn] (or) at (6,0) {};
  \node[xor gate US,draw,logic gate inputs=nn] (xor) at (9,0) {};
  \node[nand gate US,draw,logic gate inputs=nn] (nand) at (12,0) {};
  \node[below=5mm of not] {NOT};
  \node[below=5mm of and] {AND};
  \node[below=5mm of or] {OR};
  \node[below=5mm of xor] {XOR};
  \node[below=5mm of nand] {NAND};
\end{tikzpicture}
\end{center}
\diagramnote{逻辑符号表示布尔关系；具体器件还包含电源脚、阈值、传播延迟和
驱动能力。}

\topic{德摩根定律解释低有效逻辑}
\[
\lnot(A\land B)=(\lnot A)\lor(\lnot B),
\qquad
\lnot(A\lor B)=(\lnot A)\land(\lnot B).
\]

低有效信号在名称上常带 \code{\_N} 或上划线。若多个故障源都能把共享线拉低，
“任一故障发生”在正逻辑描述为 OR，在低有效线路实现中常表现为多个开漏输出
共同拉低。不能只看符号形状，必须同时看有效极性。

\topic{晶体管怎样组成反相器}
CMOS 反相器由上拉 P-MOS 与下拉 N-MOS 组成。输入低时 P-MOS 导通、N-MOS
关闭，输出趋向 \(V_{\mathrm{DD}}\)；输入高时相反，输出趋向 GND。

\begin{center}
\begin{circuitikz}
  \draw (0,4) node[vcc]{\(V_{\mathrm{DD}}\)}
    -- (0,3.4) node[pmos,anchor=S] (p) {};
  \draw (p.D) -- (0,1.8) node[nmos,anchor=D] (n) {};
  \draw (n.S) -- (0,0) node[ground]{};
  \draw (p.G) -- (-1.8,2.8);
  \draw (n.G) -- (-1.8,1.2);
  \draw (-1.8,1.2) -- (-1.8,2.8) -- (-2.8,2) node[left]{\(A\)};
  \draw (0,1.8) -- (2.0,1.8) node[right]{\(Y=\lnot A\)};
\end{circuitikz}
\end{center}

切换不是瞬间完成。输出节点具有电容，晶体管有有限导通电阻，因而边沿需要
时间。切换期间上下管还可能短暂同时部分导通。动态功耗常用近似式：

\[
P_{\mathrm{dynamic}}\approx\alpha C_L V_{\mathrm{DD}}^2 f,
\]

其中 \(\alpha\) 是翻转活动率，\(C_L\) 是负载电容，\(f\) 是切换频率。
这解释了为什么提高电压和频率会显著增加功耗，也解释了现代处理器为什么需要
时钟门控、动态调压和散热设计。

\topic{逻辑 0/1 的电气契约包含发送端和接收端}
发送端保证输出范围 \(V_{OL(max)}\)、\(V_{OH(min)}\)；接收端保证把
\(V_{IL(max)}\) 以下判为低，把 \(V_{IH(min)}\) 以上判为高。中间区域不是
“随机区”，而是手册不承诺逻辑结果的区域。

\begin{center}
\begin{tikzpicture}[x=1.05cm,y=0.9cm,font=\footnotesize\sffamily]
  \draw[->,BookMuted,thick] (0,0) -- (13.6,0) node[right]{电压};
  \fill[BookTeal!18] (0,0.25) rectangle (3.2,1.1);
  \fill[BookBlue!12] (4.2,0.25) rectangle (8.4,1.1);
  \fill[BookAmber!18] (9.5,0.25) rectangle (13.0,1.1);
  \node at (1.6,0.68) {发送器保证低输出};
  \node at (6.3,0.68) {不确定区};
  \node at (11.25,0.68) {发送器保证高输出};
  \foreach \x/\label/\colorname in {
    3.2/\(V_{OL(max)}\)/BookTeal,
    4.2/\(V_{IL(max)}\)/BookBlue,
    8.4/\(V_{IH(min)}\)/BookBlue,
    9.5/\(V_{OH(min)}\)/BookAmber} {
    \draw[\colorname,thick] (\x,0.1) -- (\x,1.45);
    \node[above,rotate=25,anchor=west,\colorname] at (\x,1.42) {\label};
  }
  \draw[<->,BookRed,thick] (3.2,-0.45) -- (4.2,-0.45)
    node[midway,below]{低噪声裕量};
  \draw[<->,BookRed,thick] (8.4,-0.45) -- (9.5,-0.45)
    node[midway,below]{高噪声裕量};
\end{tikzpicture}
\end{center}
\diagramnote{只有发送端最坏输出与接收端最坏阈值共同满足时，接口才有正噪声
裕量。图中位置是概念关系，不代表某种具体逻辑族的电压数值。}

\begin{workedexample}
发送器保证 \(V_{OL(max)}=\SI{0.4}{V}\)、
\(V_{OH(min)}=\SI{2.4}{V}\)；接收器保证
\(V_{IL(max)}=\SI{0.8}{V}\)、
\(V_{IH(min)}=\SI{2.0}{V}\)。则

\[
NM_L=0.8-0.4=\SI{0.4}{V},
\qquad
NM_H=2.4-2.0=\SI{0.4}{V}.
\]

静态上兼容。但仍需检查供电范围、输入绝对最大值、输出电流、边沿和时序；
正噪声裕量不能单独证明接口完整。
\end{workedexample}

\topic{推挽输出主动拉高也主动拉低}
CMOS 推挽输出用上拉管和下拉管分别驱动高低电平。它能提供较强的上升和下降
能力，但两个推挽输出若一个拉高、一个拉低，会形成争用大电流，不能直接并联。

\begin{center}
\begin{circuitikz}
  \draw (0,4.4) node[vcc]{\(V_{\mathrm{DD}}\)}
    -- (0,3.8) node[pmos,anchor=S] (p) {};
  \draw (p.D) -- (0,2.1) node[circ] (out) {};
  \draw (out) -- (0,1.5) node[nmos,anchor=D] (n) {};
  \draw (n.S) -- (0,0) node[ground]{};
  \draw (out) -- (2.1,2.1) node[right]{OUT};
  \draw[<-,BookBlue,thick] (p.G) -- (-2,3.1) node[left]{控制};
  \draw[<-,BookBlue,thick] (n.G) -- (-2,1.0) node[left]{互补控制};
  \node[align=center] at (0,-0.65) {推挽：主动上拉/下拉};

  \begin{scope}[xshift=7cm]
    \draw (0,4.4) node[vcc]{\(V_{\mathrm{DD}}\)}
      to[R,l={上拉 \(R_P\)}] (0,2.1) node[circ] (odout) {};
    \draw (odout) -- (0,1.5) node[nmos,anchor=D] (odn) {};
    \draw (odn.S) -- (0,0) node[ground]{};
    \draw (odout) -- (2.1,2.1) node[right]{BUS};
    \draw[<-,BookTeal,thick] (odn.G) -- (-2,1.0) node[left]{控制};
    \node[align=center] at (0,-0.65) {开漏：只主动拉低，\\释放后由电阻拉高};
  \end{scope}
\end{circuitikz}
\end{center}

\topic{开漏允许多个器件共享一根“任意一方可拉低”的线}
多个开漏输出关闭时都呈高阻，由公共电阻把线拉高；任一输出导通都把线拉低。
这适合低有效中断、I\(^2\)C 等接口。上升沿由上拉电阻和总线电容决定：

\[
\tau=R_P C_{\mathrm{bus}}.
\]

上拉太大时上升慢，太小时低电平电流增大。共享能力不是免费的。

\begin{workedexample}
\(R_P=\SI{4.7}{k\ohm}\)、总线电容 \(C=\SI{200}{pF}\)：

\[
\tau=4.7\times10^3\times200\times10^{-12}
=\SI{0.94}{\micro s}.
\]

达到 \(70\%\) 电源电压约需 \(1.204\tau\approx\SI{1.13}{\micro s}\)。若协议
要求更短上升时间，应减小电阻、降低电容、使用有源上拉或降低速率，同时检查
器件允许的灌电流。
\end{workedexample}

\topic{三态输出增加高阻态，用于分时共享}
三态输出具有 0、1 和 \(Z\)（高阻）三种电气状态。\(Z\) 不是第三个逻辑值，
而是输出暂时不驱动总线。仲裁必须确保同一时刻至多一个推挽发送器使能；全部
高阻时总线仍可能悬空，需要保持器或上下拉。

\begin{osgridlongtable}{L{0.17\linewidth}L{0.25\linewidth}L{0.28\linewidth}L{0.20\linewidth}}
输出结构 & 可主动拉高 & 可主动拉低 & 典型共享方式 \\ \hline
\endfirsthead
输出结构 & 可主动拉高 & 可主动拉低 & 典型共享方式 \\ \hline
\endhead
推挽 & 是 & 是 & 点对点；不可无仲裁并联 \\ \hline
开漏/开集 & 否 & 是 & 多设备共同拉低，外部上拉 \\ \hline
三态 & 使能时是 & 使能时是 & 通过使能/仲裁分时共享 \\ \hline
输入 & 否 & 否 & 只观察，不应驱动网络 \\ \hline
\end{osgridlongtable}

\topic{扇出由电流能力、输入电容和时序共同限制}
\term{扇出}{fan-out}是一个输出可可靠驱动的输入数量。静态时要满足

\[
N I_{IH}\le I_{OH},
\qquad
N I_{IL}\le I_{OL}.
\]

CMOS 输入静态电流很小，但每个输入和走线增加电容，数量越多边沿越慢：

\[
t_r\propto R_{\mathrm{driver}}C_{\mathrm{load}}.
\]

所以“输入几乎不耗电”不等于一个 GPIO 能无限分叉。

\begin{workedexample}
输出允许灌电流 \(\SI{8}{mA}\)，每个接收器低电平输入电流最坏
\(\SI{0.4}{mA}\)，静态电流给出的扇出上限为

\[
N\le\frac{8}{0.4}=20.
\]

若时序分析发现驱动 12 个输入时上升沿已经超过协议上限，则实际扇出必须小于
12。静态和动态两个约束取更严格者。
\end{workedexample}

\topic{CMOS NAND 与 NOR 怎样从串并联晶体管得到}
CMOS 门由互补上拉网络 PUN 和下拉网络 PDN 组成。对 2-input NAND：
N-MOS 串联，只有 \(A=B=1\) 才能下拉；P-MOS 并联，任一输入为 0 就能上拉。

\begin{center}
\begin{circuitikz}
  \draw (0,5.5) node[vcc]{\(V_{\mathrm{DD}}\)};
  \draw (-1.1,5.0) node[pmos,anchor=S] (p1) {};
  \draw (1.1,5.0) node[pmos,anchor=S] (p2) {};
  \draw (0,5.5) -| (p1.S);
  \draw (0,5.5) -| (p2.S);
  \draw (p1.D) |- (0,3.8) node[circ] (y) {};
  \draw (p2.D) |- (y);
  \draw (y) -- (3.0,3.8) node[right]{\(Y\)};
  \draw (0,3.8) -- (0,3.2) node[nmos,anchor=D] (n1) {};
  \draw (n1.S) -- (0,1.7) node[nmos,anchor=D] (n2) {};
  \draw (n2.S) -- (0,0) node[ground]{};
  \draw (p1.G) -- (-2.2,4.2) node[left]{\(A\)};
  \draw (n1.G) -- (-2.2,2.6) node[left]{\(A\)};
  \draw (p2.G) -- (2.2,4.2) node[right]{\(B\)};
  \draw (n2.G) -- (2.2,1.1) node[right]{\(B\)};
  \node[align=center] at (0,-0.7) {NAND：上拉并联、下拉串联};

  \begin{scope}[xshift=8cm]
    \node[os block,minimum width=24mm] (nand) at (0,4.0) {NAND};
    \node[os block,minimum width=24mm] (inv) at (0,1.9) {反相器};
    \node[left=15mm of nand,yshift=3mm] (a) {\(A\)};
    \node[left=15mm of nand,yshift=-3mm] (b) {\(B\)};
    \node[right=15mm of inv] (andout) {\(A\land B\)};
    \draw[os arrow] (a) -- ([yshift=3mm]nand.west);
    \draw[os arrow] (b) -- ([yshift=-3mm]nand.west);
    \draw[os arrow] (nand) -- node[right]{取反} (inv);
    \draw[os arrow] (inv) -- (andout);
    \node[align=center] at (0,0.3)
      {NAND 后再反相得到 AND；\\功能完备描述可构造性};
  \end{scope}
\end{circuitikz}
\end{center}

NOR 使用并联 N-MOS 下拉和串联 P-MOS 上拉。晶体管级图让德摩根定律、有效
极性和导电路径连接起来，而不仅是一张真值表。

\topic{电压域不同必须有明确的电平转换策略}
\begin{osgridlongtable}{L{0.23\linewidth}L{0.31\linewidth}L{0.36\linewidth}}
情况 & 可能方案 & 必须核对 \\ \hline
\endfirsthead
情况 & 可能方案 & 必须核对 \\ \hline
\endhead
低压推挽驱动高压输入 & 输入阈值兼容时可直连，否则缓冲/转换 & \(V_{IH}\)、输入容限、边沿 \\ \hline
高压推挽驱动低压输入 & 分压、限流钳位或专用转换器 & 绝对最大、速度、掉电反灌 \\ \hline
开漏双向总线 & 双电源 MOSFET 或专用双向转换器 & 上拉、电容、方向与关机状态 \\ \hline
高速差分接口 & 协议专用 PHY/收发器 & 共模、终端、阻抗、参考时钟 \\ \hline
\end{osgridlongtable}

\begin{failurebox}
输入端“能读出正确 0/1”不等于连接安全。高压可能经输入保护二极管向未上电
电源域反灌；慢边沿可能增加短路电流或重复越过阈值；多个推挽输出可能争用。
逻辑功能、电气安全和时序必须分别验收。
\end{failurebox}


\section{布尔代数与组合逻辑}

\term{组合逻辑}{combinational logic}没有有意保存的历史状态。在传播延迟
过去后，输出由当前输入唯一决定。

\topic{半加器与全加器}
半加器把两个 1-bit 数相加：

\[
S=A\oplus B,\qquad C=A\land B.
\]

\begin{center}
\begin{osgridtabular}{cc|cc}
\(A\) & \(B\) & 和 \(S\) & 进位 \(C\) \\ \hline
0 & 0 & 0 & 0\\ \hline
0 & 1 & 1 & 0\\ \hline
1 & 0 & 1 & 0\\ \hline
1 & 1 & 0 & 1\\ \hline
\end{osgridtabular}
\end{center}

全加器还接收低位进位 \(C_{\mathrm{in}}\)：

\[
S=A\oplus B\oplus C_{\mathrm{in}},
\]
\[
C_{\mathrm{out}}
=(A\land B)\lor
\left(C_{\mathrm{in}}\land(A\oplus B)\right).
\]

\begin{pdfdiagram}
  \node[os block] (fa0) {bit 0\\全加器};
  \node[os block,right=of fa0] (fa1) {bit 1\\全加器};
  \node[os block,right=of fa1] (fa2) {\(\cdots\)};
  \node[os block,right=of fa2] (fan) {最高 bit\\全加器};
  \draw[os arrow] (fa0) -- node[above]{进位} (fa1);
  \draw[os arrow] (fa1) -- node[above]{进位} (fa2);
  \draw[os arrow] (fa2) -- node[above]{进位} (fan);
\end{pdfdiagram}
\diagramnote{最简单的串行进位加法器必须等待进位逐级传播。更快的加法器用
额外逻辑并行预测进位，以面积和布线换取延迟。}

\begin{workedexample}
计算四位 \(1011_2+0110_2\)。从最低位开始：

\[
\begin{array}{c@{\quad}c@{\quad}c@{\quad}c}
1+0=1, &
1+1=0\ \text{进 1}, &
0+1+1=0\ \text{进 1}, &
1+0+1=0\ \text{进 1}.
\end{array}
\]

得到五位结果 \(1\,0001_2=17\)。若结果被限制在 4 bit，只保留
\(0001_2\)，并用 carry flag 记录无符号进位。位宽是计算语义的一部分。
\end{workedexample}

\topic{多路选择器：由控制位选择数据来源}
2 选 1 多路选择器满足

\[
Y=(\lnot S\land D_0)\lor(S\land D_1).
\]

\begin{pdfdiagram}
  \node[trapezium,trapezium left angle=70,trapezium right angle=110,
        draw=BookBlue,fill=BookCodeBg,minimum width=18mm,minimum height=20mm] (mux) {MUX};
  \node[left=16mm of mux,yshift=5mm] (d0) {\(D_0\)};
  \node[left=16mm of mux,yshift=-5mm] (d1) {\(D_1\)};
  \node[below=10mm of mux] (s) {\(S\)};
  \node[right=16mm of mux] (y) {\(Y\)};
  \draw[os arrow] (d0) -- (mux.west |- d0);
  \draw[os arrow] (d1) -- (mux.west |- d1);
  \draw[os arrow] (s) -- (mux.south);
  \draw[os arrow] (mux) -- (y);
\end{pdfdiagram}

CPU 数据通路中，多路选择器决定 ALU 输入来自寄存器、立即数、程序计数器还是
转发结果；控制器产生选择信号。

\topic{译码器：把编号变成唯一选择}
一个 \(n\)-to-\(2^n\) 译码器把 \(n\) 位输入展开为最多 \(2^n\) 条选择线。
地址译码就是更大规模的选择：比较地址的某些高位，决定当前事务送往 RAM、
ROM 还是设备。

\begin{workedexample}
若一个芯片只在地址 \(\texttt{0xF000}\)--\(\texttt{0xFFFF}\) 响应，窗口大小为

\[
\texttt{0x10000}-\texttt{0xF000}=\texttt{0x1000}=4096\ \mathrm{B}.
\]

在 16-bit 地址中，这等价于最高 4 bit 必须为 \(1111_2\)，低 12 bit 选择
窗口内的 4096 个字节。地址比较负责选芯，低位负责片内寻址。
\end{workedexample}

\topic{布尔代数把真值表变成可化简的表达式}
常用恒等律包括：

\[
A\land1=A,\quad A\lor0=A,\qquad
A\land0=0,\quad A\lor1=1,
\]
\[
A\land A=A,\quad A\lor A=A,\qquad
A\land\lnot A=0,\quad A\lor\lnot A=1,
\]
\[
A\land(B\lor C)=(A\land B)\lor(A\land C).
\]

逻辑化简可以减少门数量、减少门级，或适配已有标准单元。但表达式更短不总
等于物理实现更快；布局、扇出与门类型同样影响延迟。

\topic{从真值表写出“与项之和”}
对输出为 1 的每一行，写一个只在该行成立的与项，再把这些项 OR。例如多数
表决函数 \(M(A,B,C)\) 在至少两个输入为 1 时输出 1：

\[
M=(A\land B\land\lnot C)
\lor(A\land\lnot B\land C)
\lor(\lnot A\land B\land C)
\lor(A\land B\land C).
\]

化简后：

\[
M=(A\land B)\lor(A\land C)\lor(B\land C).
\]

这也正是全加器进位输出的一种表达形式。

\begin{center}
\begin{osgridtabular}{ccc|c}
\(A\) & \(B\) & \(C\) & 多数 \(M\) \\ \hline
0&0&0&0\\ \hline
0&0&1&0\\ \hline
0&1&0&0\\ \hline
0&1&1&1\\ \hline
1&0&0&0\\ \hline
1&0&1&1\\ \hline
1&1&0&1\\ \hline
1&1&1&1\\ \hline
\end{osgridtabular}
\end{center}

\topic{卡诺图用相邻格合并只差一个 bit 的项}
卡诺图按 Gray code 排列，让相邻格只有一个输入位变化。下图函数在 \(B=1\)
的两行都为 1，因此 \(A\) 可被消去：

\begin{center}
\begin{tikzpicture}[font=\footnotesize\sffamily]
  \draw (0,0) grid[xstep=2,ystep=1] (4,2);
  \node at (1,2.35) {\(B=0\)};
  \node at (3,2.35) {\(B=1\)};
  \node[left] at (0,1.5) {\(A=0\)};
  \node[left] at (0,0.5) {\(A=1\)};
  \node at (1,1.5) {0};
  \node at (3,1.5) {1};
  \node at (1,0.5) {0};
  \node at (3,0.5) {1};
  \draw[BookBlue,very thick,rounded corners=8pt]
    (2.12,0.12) rectangle (3.88,1.88);
  \node[BookBlue,right] at (4.1,1.0) {合并后 \(Y=B\)};
\end{tikzpicture}
\end{center}

卡诺图适合少量变量的纸笔化简；大规模逻辑由综合工具完成，但工程师仍需理解
don’t-care、冒险和时序约束。

\topic{半加器和全加器应该能从门级逐线读出}
\begin{center}
\begin{tikzpicture}[circuit logic US,font=\footnotesize\sffamily,>=Latex]
  \node (a) at (0,2.2) {\(A\)};
  \node (b) at (0,0.8) {\(B\)};
  \node[xor gate US,draw,logic gate inputs=nn] (xor) at (4,2.0) {};
  \node[and gate US,draw,logic gate inputs=nn] (and) at (4,0.4) {};
  \draw[->] (a) -- ++(1,0) |- (xor.input 1);
  \draw[->] (b) -- ++(1.4,0) |- (xor.input 2);
  \draw[->] (a) -- ++(1.0,0) |- (and.input 1);
  \draw[->] (b) -- ++(1.4,0) |- (and.input 2);
  \draw[->] (xor.output) -- ++(1.4,0) node[right]{和 \(S\)};
  \draw[->] (and.output) -- ++(1.4,0) node[right]{进位 \(C\)};
  \node[align=center] at (3,-0.7)
    {半加器：XOR 产生本位和，AND 产生向高位的进位};

  \begin{scope}[xshift=8cm]
    \node[os block,minimum width=19mm] (ha1) at (0,2.1) {半加器 1};
    \node[os block,minimum width=19mm] (ha2) at (3.0,2.1) {半加器 2};
    \node[or gate US,draw,logic gate inputs=nn] (or) at (4.7,0.4) {};
    \node[left=12mm of ha1,yshift=3mm] (fa) {\(A\)};
    \node[left=12mm of ha1,yshift=-3mm] (fb) {\(B\)};
    \node[above=8mm of ha2] (cin) {\(C_{\mathrm{in}}\)};
    \draw[os arrow] (fa) -- ([yshift=3mm]ha1.west);
    \draw[os arrow] (fb) -- ([yshift=-3mm]ha1.west);
    \draw[os arrow] (ha1) -- node[above]{部分和} (ha2);
    \draw[os arrow] (cin) -- (ha2);
    \draw[os arrow] (ha2) -- ++(1.4,0) node[right]{\(S\)};
    \draw[os arrow] (ha1.south) |- (or.input 1);
    \draw[os arrow] (ha2.south) |- (or.input 2);
    \draw[os arrow] (or.output) -- ++(1.0,0) node[right]{\(C_{\mathrm{out}}\)};
  \end{scope}
\end{tikzpicture}
\end{center}
\diagramnote{全加器可由两个半加器和一个 OR 组成。实际标准单元可能采用更快
拓扑，但逻辑关系相同。}

\topic{进位链决定宽加法器延迟}
串行进位加法器中，第 \(i\) 位必须等待低位进位。定义

\[
G_i=A_i\land B_i\quad\text{（本位生成进位）},
\]
\[
P_i=A_i\oplus B_i\quad\text{（本位传播进位）},
\]

则

\[
C_{i+1}=G_i\lor(P_i\land C_i).
\]

展开多位关系可以并行预测进位，形成 carry-lookahead；树形前缀加法器进一步
降低逻辑深度。代价是更多门、布线和功耗。

\topic{carry 与 signed overflow 来自不同条件}
无符号进位关注最高位之外的 carry；补码溢出发生在两个同号操作数得到异号
结果时。对最高位进位还可写成

\[
OF=C_{n-1}\oplus C_n.
\]

\begin{workedexample}
8-bit 模式 \(\texttt{1111\,1110}\) 与 \(\texttt{0000\,0010}\)：

\begin{itemize}[leftmargin=2.2em]
  \item 无符号解释：254 大于 2；
  \item 补码解释：\(-2\) 小于 2。
\end{itemize}

bit 完全相同，比较结果不同，原因是比较器选择了不同数值解释。CPU 条件跳转
因此分别提供 unsigned 和 signed 条件组合。
\end{workedexample}

\topic{译码器、编码器和优先编码器方向相反}
\begin{osgridlongtable}{L{0.21\linewidth}L{0.30\linewidth}L{0.39\linewidth}}
模块 & 输入到输出 & 典型用途 \\ \hline
\endfirsthead
模块 & 输入到输出 & 典型用途 \\ \hline
\endhead
译码器 & \(n\) 位编号 \(\rightarrow 2^n\) 条 one-hot 线 & 寄存器选择、片选、指令类别 \\ \hline
编码器 & one-hot 输入 \(\rightarrow n\) 位编号 & 把唯一请求者变成索引 \\ \hline
优先编码器 & 多个请求 \(\rightarrow\) 最高优先级编号与 valid & 中断、仲裁和前导零检测 \\ \hline
\end{osgridlongtable}

\begin{center}
\begin{tikzpicture}[font=\footnotesize\sffamily,>=Latex]
  \node[draw=BookBlue,fill=BookCodeBg,trapezium,trapezium left angle=70,
        trapezium right angle=110,minimum width=25mm,minimum height=26mm,
        align=center]
        (dec) at (0,2) {2-to-4\\译码器};
  \foreach \shift/\name in {-9/Y0,-3/Y1,3/Y2,9/Y3}
    \draw[->] ([yshift=\shift mm]dec.east) -- ++(1.25,0) node[right]{\name};
  \draw[->] (-1.5,1.6) node[left]{\(A_0\)} -- ([yshift=-4mm]dec.west);
  \draw[->] (-1.5,2.4) node[left]{\(A_1\)} -- ([yshift=4mm]dec.west);

  \node[draw=BookTeal,fill=BookCodeBg,trapezium,trapezium left angle=110,
        trapezium right angle=70,minimum width=28mm,minimum height=28mm,
        align=center]
        (enc) at (6.8,2) {4-to-2\\优先编码器};
  \foreach \shift/\name in {-9/R0,-3/R1,3/R2,9/R3}
    \draw[->] ([yshift=\shift mm]enc.west)+(-1.25,0) node[left]{\name}
      -- ([yshift=\shift mm]enc.west);
  \draw[->] ([yshift=-4mm]enc.east) -- ++(1.4,0) node[right]{INDEX};
  \draw[->] ([yshift=4mm]enc.east) -- ++(1.4,0) node[right]{VALID};
\end{tikzpicture}
\end{center}

\topic{桶形移位器用多级 MUX 在固定级数内移动多位}
逐次移动 \(k\) 次会让延迟随 \(k\) 增长。8-bit 桶形移位器可使用 3 级
多路选择器，分别选择移动 0/1、0/2、0/4 位，组合得到任意距离。

\begin{pdfdiagram}[node distance=7mm and 11mm]
  \node[os state,minimum width=22mm] (input) {8-bit 输入};
  \node[os block,right=of input,minimum width=24mm] (s1) {第 0 级\\移 0/1};
  \node[os block,right=of s1,minimum width=24mm] (s2) {第 1 级\\移 0/2};
  \node[os block,right=of s2,minimum width=24mm] (s4) {第 2 级\\移 0/4};
  \node[os state,right=of s4,minimum width=22mm] (output) {8-bit 输出};
  \draw[os bus] (input) -- (s1);
  \draw[os bus] (s1) -- (s2);
  \draw[os bus] (s2) -- (s4);
  \draw[os bus] (s4) -- (output);
  \node[below=8mm of s2,align=center]
    {控制位 \(\texttt{shift[2:0]}\) 分别控制三层};
\end{pdfdiagram}

\topic{组合逻辑必须定义非法输入和输出有效条件}
如果译码器输入含越界组合，输出是全不选、选默认目标还是报告错误？如果除法器
除数为 0，结果和异常怎样表达？如果多个请求同时有效，优先级如何决定？高质量
接口要同时定义正常路径、默认值和错误路径。

\begin{failurebox}
真值表只列 0/1 时容易隐藏现实问题：输入可能在阈值间、正在切换、违反时序，
或来自不同电压/时钟域。组合逻辑正确性需要布尔关系、电气范围和稳定时间三份
条件同时成立。
\end{failurebox}


\section{传播延迟、反馈与状态存储}

一个门从输入改变到输出稳定需要传播延迟 \(t_{pd}\)。不同路径延迟不同，即使
布尔表达式最终正确，中间也可能产生短脉冲，称为 glitch 或 hazard。

\begin{center}
\begin{tikzpicture}[x=1cm,y=0.65cm,font=\footnotesize\sffamily]
  \draw[->] (0,0) -- (12,0) node[right]{时间};
  \foreach \y/\name in {3/A,2/B,1/Y} {
    \draw[BookRule] (0,\y) -- (12,\y);
    \node[left] at (0,\y) {\name};
  }
  \draw[BookBlue,thick] (0,3) -- (2,3) -- (2,3.45) -- (8,3.45) -- (8,3) -- (12,3);
  \draw[BookTeal,thick] (0,2) -- (5,2) -- (5,2.45) -- (10,2.45) -- (10,2);
  \draw[BookRed,thick] (0,1.45) -- (5.4,1.45) -- (5.4,1)
    -- (5.9,1) -- (5.9,1.45) -- (12,1.45);
  \draw[<->] (5,0.45) -- (5.9,0.45) node[midway,below]{短暂毛刺};
\end{tikzpicture}
\end{center}

同步设计通常在时钟边沿采样，让组合逻辑有一个周期减去时序裕量的时间稳定。
异步复位、时钟切换和外部输入仍需专门处理，不能假定所有毛刺都会自动消失。

\topic{存储状态：从反馈到锁存器}
若电路输出反馈到输入，就可能形成多个稳定状态。交叉耦合门组成的 SR latch
能够保存一个 bit，但某些输入组合非法或具有歧义。D latch 用额外逻辑把数据
和使能组织成更安全接口；D flip-flop 只在规定时钟边沿采样。

\begin{pdfdiagram}
  \node[os state,minimum width=26mm] (source) {数据源\\在边沿前稳定 \(D\)};
  \node[draw=BookBlue,fill=BookCodeBg,minimum width=31mm,minimum height=24mm,
        align=center,right=22mm of source] (ff) {D 触发器\\\(D\quad\triangleright\quad Q\)};
  \node[os state,minimum width=28mm,right=22mm of ff] (q) {保存输出 \(Q\)\\边沿后延迟更新};
  \node[os block,below=11mm of ff,minimum width=27mm] (clk) {时钟 CLK\\有效边沿};
  \draw[os bus] (source) -- node[above]{采样输入} (ff);
  \draw[os arrow] (clk) -- (ff);
  \draw[os bus] (ff) -- node[above]{\(t_{cq}\) 后稳定} (q);
  \draw[decorate,decoration={brace,mirror,amplitude=4pt},BookTeal]
    ($(q.south west)+(0,-5mm)$) -- ($(q.south east)+(0,-5mm)$)
    node[midway,below=5pt,BookTeal]{保持到下一有效边沿};
\end{pdfdiagram}

\topic{建立时间、保持时间与 clock-to-Q}
数据必须在采样边沿前稳定至少 \(t_{setup}\)，边沿后继续稳定至少
\(t_{hold}\)。采样后输出还需 \(t_{cq}\) 才稳定。两个寄存器之间的组合路径
要满足近似约束：

\[
T_{\mathrm{clk}}
\ge t_{cq}+t_{\mathrm{comb,max}}+t_{setup}+t_{\mathrm{skew}}.
\]

\begin{center}
\begin{tikzpicture}[x=1cm,y=0.7cm,font=\footnotesize\sffamily]
  \draw[->] (0,0) -- (12,0) node[right]{时间};
  \node[left] at (0,3) {CLK};
  \node[left] at (0,1.6) {D};
  \draw[BookBlue,thick] (0,3) -- (2,3) -- (2,3.5) -- (5,3.5)
    -- (5,3) -- (8,3) -- (8,3.5) -- (11,3.5) -- (11,3);
  \draw[dashed,BookRed] (8,0.5) -- (8,4);
  \draw[BookTeal,thick] (0,1.6) -- (6.3,1.6) -- (6.3,2.1) -- (9.2,2.1) -- (9.2,1.6) -- (12,1.6);
  \draw[decorate,decoration={brace,mirror,amplitude=4pt}]
    (6.3,0.9) -- (8,0.9) node[midway,below=5pt]{建立窗口};
  \draw[decorate,decoration={brace,mirror,amplitude=4pt}]
    (8,0.9) -- (9.2,0.9) node[midway,below=5pt]{保持窗口};
  \node[above] at (8,4) {采样边沿};
\end{tikzpicture}
\end{center}

\topic{亚稳态不是“随机读错一个 bit”这么简单}
若异步输入恰在采样窗口变化，触发器内部可能暂时停留在不满足正常 0/1 时序的
状态，称为 metastability。它最终通常会解析到 0 或 1，但解析时间无法对单次
事件作确定保证。常见做法是在目标时钟域串联同步触发器，把亚稳态传播到业务
逻辑的概率降得足够低；多 bit 总线则需要握手、FIFO 或编码协议，不能让每一位
独立同步后就假定组合仍一致。

\section{触发器、时钟与跨时钟域}

下面的低有效 NAND 型 SR latch 用
\(\overline{S}\) 置位、\(\overline{R}\) 复位。两输入都为 1 时保持；把两者
同时拉低会让两个输出都为 1，释放顺序又可能决定最终状态，因此该组合不允许。

\begin{center}
\begin{tikzpicture}[circuit logic US,font=\footnotesize\sffamily,>=Latex]
  \node[nand gate US,draw,logic gate inputs=nn] (top) at (4,3.1) {};
  \node[nand gate US,draw,logic gate inputs=nn] (bottom) at (4,0.8) {};
  \node[left=22mm of top.input 1] (sn) {\(\overline{S}\)};
  \node[left=22mm of bottom.input 2] (rn) {\(\overline{R}\)};
  \draw[->] (sn) -- (top.input 1);
  \draw[->] (rn) -- (bottom.input 2);
  \draw[->,BookBlue,thick] (top.output) -- ++(1.7,0)
    node[right] {\(Q\)};
  \draw[->,BookTeal,thick] (bottom.output) -- ++(1.7,0)
    node[right] {\(\overline{Q}\)};
  \draw[->,BookTeal,thick] (top.output) -- ++(0.7,0)
    |- (bottom.input 1);
  \draw[->,BookBlue,thick] (bottom.output) -- ++(0.45,0)
    |- (top.input 2);
  \node[align=center] at (4,-0.4)
    {交叉反馈让输出在输入释放后仍保持};

  \node[draw=BookRule,fill=white,minimum width=43mm,minimum height=30mm,
        align=left] at (11,2.0)
    {\(\overline{S}\ \overline{R}\quad Q_{\mathrm{next}}\)\\
     \(1\quad1\qquad Q_{\mathrm{current}}\)\\
     \(0\quad1\qquad1\)（置位）\\
     \(1\quad0\qquad0\)（复位）\\
     \(0\quad0\qquad\)非法};
\end{tikzpicture}
\end{center}
\diagramnote{反馈连线分别从 \(Q\) 回到底部门、从 \(\overline{Q}\) 回到顶部门；
输出线与反馈线分开绘制，避免把交叉点误读为短接。}

\topic{D latch 用使能窗口控制“何时透明”}
D latch 在 EN 有效期间，输出会在传播延迟后跟随 D；EN 失效后保持最后状态。
因此它是电平敏感器件，不是只在一个瞬间采样。

\begin{center}
\begin{tikzpicture}[x=0.9cm,y=0.72cm,font=\footnotesize\sffamily]
  \draw[->,BookMuted] (0,0) -- (14,0) node[right]{时间};
  \foreach \y/\name in {4/EN,2.6/D,1.2/Q} {
    \draw[BookRule] (0,\y) -- (13.5,\y);
    \node[left] at (0,\y) {\name};
  }
  \fill[BookBlue!10] (2,0.45) rectangle (8,4.8);
  \node[BookBlue] at (5,4.55) {透明窗口};
  \draw[BookBlue,thick] (0,4) -- (2,4) -- (2,4.45) -- (8,4.45) -- (8,4) -- (13.5,4);
  \draw[BookTeal,thick] (0,2.6) -- (3,2.6) -- (3,3.05) -- (5.3,3.05)
    -- (5.3,2.6) -- (7,2.6) -- (7,3.05) -- (10.5,3.05) -- (10.5,2.6) -- (13.5,2.6);
  \draw[BookAmber,thick] (0,1.2) -- (3.25,1.2) -- (3.25,1.65)
    -- (5.55,1.65) -- (5.55,1.2) -- (7.25,1.2) -- (7.25,1.65) -- (13.5,1.65);
  \draw[dashed,BookRed] (8,0.5) -- (8,4.8);
  \node[BookRed,below] at (8,0.5) {关闭后保持};
\end{tikzpicture}
\end{center}

\topic{边沿触发器把采样窗口压缩到时钟边沿附近}
实际边沿触发 DFF 可由主从锁存器或专用晶体管结构实现。接口只承诺：D 在
建立/保持窗口内稳定时，Q 在 clock-to-Q 延迟后更新为采样值，并保持到下一
有效边沿。

\begin{pdfdiagram}[node distance=8mm and 12mm]
  \node[os state,minimum width=25mm] (source) {源寄存器\\边沿后产生 \(Q_1\)};
  \node[os block,right=of source,minimum width=31mm] (logic) {组合逻辑\\最大延迟 \(t_{\mathrm{comb,max}}\)};
  \node[os state,right=of logic,minimum width=25mm] (dest) {目标寄存器\\采样 \(D_2\)};
  \draw[os bus] (source) -- node[above,font=\scriptsize]{\(t_{cq}\)} (logic);
  \draw[os bus] (logic) -- node[above,font=\scriptsize]{建立前到达} (dest);
  \draw[BookBlue,thick] ($(source.south)+(0,-8mm)$) -- ($(dest.south)+(0,-8mm)$);
  \draw[BookBlue,thick] ($(source.south)+(0,-10mm)$) -- ($(source.south)+(0,-6mm)$);
  \draw[BookBlue,thick] ($(dest.south)+(0,-10mm)$) -- ($(dest.south)+(0,-6mm)$);
  \node[below=11mm of logic] {一个时钟周期预算 \(T_{\mathrm{clk}}\)};
\end{pdfdiagram}

建立约束防止数据到得太晚：

\[
T_{\mathrm{clk}}\ge
t_{cq,\max}+t_{\mathrm{comb,max}}+t_{setup}+t_{\mathrm{skew}}.
\]

保持约束防止新数据到得太快：

\[
t_{cq,\min}+t_{\mathrm{comb,min}}
\ge t_{hold}+t_{\mathrm{skew}}.
\]

降低时钟频率能放宽建立约束，却不自动修复保持违例，因为保持问题发生在同一
边沿附近。

\topic{时钟偏斜、抖动和占空比是三种不同误差}
\begin{osgridlongtable}{L{0.20\linewidth}L{0.31\linewidth}L{0.39\linewidth}}
概念 & 含义 & 可能后果 \\ \hline
\endfirsthead
概念 & 含义 & 可能后果 \\ \hline
\endhead
clock skew & 同一名义边沿到不同寄存器的时间差 & 吃掉建立或保持裕量 \\ \hline
jitter & 同一位置的边沿相对理想时刻抖动 & 周期预算变化、采样误差 \\ \hline
duty cycle distortion & 高低电平宽度偏离目标比例 & 半周期路径或 PLL/门控违例 \\ \hline
clock uncertainty & 设计中合并考虑的时序不确定量 & 时序分析中预留保守裕量 \\ \hline
\end{osgridlongtable}

时钟不是普通数据线。随意用组合门对时钟做 AND/OR 可能产生窄脉冲或额外
偏斜；时钟门控通常使用专用单元，在安全相位锁存 enable。

\topic{两级同步器降低单 bit 亚稳态传播概率}
\begin{center}
\begin{tikzpicture}[font=\footnotesize\sffamily,>=Latex]
  \node[os hardware] (async) at (0,2) {异步输入};
  \node[os state] (ff1) at (4,2) {DFF 1\\可能亚稳};
  \node[os state] (ff2) at (8,2) {DFF 2\\业务可用};
  \node[os block] (logic) at (12,2) {目标域逻辑};
  \draw[os arrow] (async) -- (ff1);
  \draw[os arrow] (ff1) -- node[above]{一周期解析时间} (ff2);
  \draw[os arrow] (ff2) -- (logic);
  \draw[BookBlue,thick] (3,0.5) -- (9,0.5);
  \node[below] at (6,0.5) {两级都由目标时钟 CLK\(_B\) 采样};
  \draw[os arrow] (4,0.5) -- (ff1.south);
  \draw[os arrow] (8,0.5) -- (ff2.south);
\end{tikzpicture}
\end{center}

常见概率模型把平均失效间隔近似写成

\[
MTBF\propto
\frac{e^{T_{\mathrm{resolve}}/\tau_m}}
{f_{\mathrm{clk}}f_{\mathrm{data}}},
\]

其中参数依工艺和单元。增加解析时间可指数级改善 MTBF，却无法给单次事件提供
绝对保证。

\topic{窄脉冲可能在目标域两个边沿之间完全消失}
源域的单周期脉冲若比目标时钟周期短，两级同步器可能一次也采不到。解决方法
包括把脉冲展宽、转换为 toggle、使用 request/acknowledge 握手，或用异步 FIFO
传递事件和数据。

\topic{多 bit 总线不能逐位各接两个触发器}
每一位独立同步可能在不同周期稳定，接收方看到源端从未产生过的混合值。常见
方案：

\begin{itemize}[leftmargin=2.2em]
  \item 数据保持稳定，同时同步 valid/ack 握手；
  \item 单调计数器采用 Gray code，让相邻状态只变一位；
  \item 连续流数据使用带 Gray 指针的异步 FIFO；
  \item 复位、时钟和电源域跨越使用专用 CDC/RDC 方法检查。
\end{itemize}

\topic{机械按键需要去抖，异步输入需要同步}
机械触点闭合时会在几毫秒内多次弹跳。按下一次可能形成一串高低边沿。RC 与
施密特触发器可以在硬件侧滤波；软件也可在首次变化后等待稳定窗口，再确认
状态。去抖解决重复边沿，同步器解决时钟域采样，二者不是同一问题。

\begin{center}
\begin{tikzpicture}[x=0.85cm,y=0.75cm,font=\footnotesize\sffamily]
  \draw[->,BookMuted] (0,0) -- (14,0) node[right]{时间};
  \node[left] at (0,3.2) {原始触点};
  \node[left] at (0,1.4) {去抖结果};
  \draw[BookRed,thick] (0,3.2) -- (2,3.2)
    -- (2,3.7) -- (2.4,3.7) -- (2.4,3.2)
    -- (2.8,3.2) -- (2.8,3.7) -- (3.1,3.7)
    -- (3.1,3.2) -- (3.5,3.2) -- (3.5,3.7) -- (13,3.7);
  \fill[BookRed!10] (2,0.5) rectangle (5.5,4.3);
  \node[BookRed] at (3.75,4.05) {弹跳窗口};
  \draw[BookTeal,very thick] (0,1.4) -- (5.5,1.4) -- (5.5,1.9) -- (13,1.9);
  \draw[dashed,BookTeal] (5.5,0.5) -- (5.5,4.3);
  \node[BookTeal,below] at (5.5,0.5) {确认稳定后只产生一次变化};
\end{tikzpicture}
\end{center}

\topic{复位也是跨时钟域信号}
异步复位可以在时钟未运行时立即强制状态，但释放若靠近时钟边沿，会违反
recovery/removal 要求，使同一状态机不同触发器在不同周期退出复位。常用结构
是“异步置位、同步释放”：每个时钟域有自己的两级释放链。

\begin{pdfdiagram}[node distance=7mm and 11mm]
  \node[os hardware,minimum width=27mm] (external) {外部 RESET\_N\\可异步拉低};
  \node[os state,right=of external,minimum width=25mm] (sync1) {释放同步 DFF 1\\复位时清零};
  \node[os state,right=of sync1,minimum width=25mm] (sync2) {释放同步 DFF 2\\目标域复位};
  \node[os block,right=of sync2,minimum width=26mm] (domain) {时钟域状态机\\同边沿退出};
  \draw[os arrow] (external) -- node[above,font=\scriptsize]{置位可异步} (sync1);
  \draw[os arrow] (sync1) -- (sync2);
  \draw[os arrow] (sync2) -- (domain);
  \node[below=8mm of sync2] {同步链由该域时钟驱动};
\end{pdfdiagram}

\topic{FSM 设计要有状态表、转移图和默认恢复}
以简化读事务控制器为例：

\begin{osgridlongtable}{L{0.18\linewidth}L{0.25\linewidth}L{0.28\linewidth}L{0.19\linewidth}}
当前状态 & 输入条件 & 下一状态 & 输出动作 \\ \hline
\endfirsthead
当前状态 & 输入条件 & 下一状态 & 输出动作 \\ \hline
\endhead
Idle & request=0 & Idle & ready=1 \\ \hline
Idle & request=1 & Address & 锁存地址，ready=0 \\ \hline
Address & 始终 & Wait & 发出 read valid \\ \hline
Wait & response=0 & Wait & 保持事务 \\ \hline
Wait & response=1,error=0 & Complete & 锁存返回数据 \\ \hline
Wait & response=1,error=1 & Error & 锁存错误码 \\ \hline
Complete & 始终 & Idle & done 脉冲 \\ \hline
Error & clear=1 & Idle & 清错误 \\ \hline
\end{osgridlongtable}

\begin{center}
\begin{tikzpicture}[font=\footnotesize\sffamily,>=Latex,node distance=18mm and 22mm]
  \node[os state] (idle) {Idle};
  \node[os state,right=of idle] (address) {Address};
  \node[os state,right=of address] (wait) {Wait};
  \node[os state,above right=of wait] (complete) {Complete};
  \node[os state,below right=of wait] (error) {Error};
  \draw[os arrow] (idle) -- node[above]{request} (address);
  \draw[os arrow] (address) -- node[above]{发出事务} (wait);
  \draw[os arrow] (wait) -- node[above,sloped]{response \(\land\lnot error\)} (complete);
  \draw[os arrow] (wait) -- node[below,sloped]{response \(\land error\)} (error);
  \draw[os arrow,bend left=25] (complete) to node[above]{done} (idle);
  \draw[os arrow,bend right=25] (error) to node[below]{clear} (idle);
  \draw[os arrow,loop below] (wait) to node[below]{\(\lnot response\)} (wait);
\end{tikzpicture}
\end{center}

非法状态必须有策略：回到 reset/idle、保持并报告错误，或触发安全停机。只写
“正常情况下永远不会进入”不能替代恢复设计与断言。

\begin{failurebox}
时序电路问题不能只靠重新画布尔表达式修复。建立/保持违例、亚稳态、CDC、
复位释放和时钟门控都与时间相关。仿真中理想 0 延迟通过，不代表真实硅片在
PVT（工艺、电压、温度）角落满足时序。
\end{failurebox}


\section{寄存器、状态机与算术逻辑单元}

把多个触发器并列可保存一个固定宽度的字，称为寄存器。寄存器前放多路选择器，
可以选择保持旧值、装载新值、加一或清零。

\begin{pdfdiagram}
  \node[os block] (next) {组合逻辑\\计算 next state};
  \node[os state,right=25mm of next] (reg) {状态寄存器\\在时钟边沿提交};
  \node[os block,right=25mm of reg] (output) {输出逻辑\\驱动数据通路};
  \draw[os arrow] (next) -- (reg);
  \draw[os arrow] (reg) -- (output);
  \draw[os arrow,bend left=35] (reg.north) to node[above]{current state} (next.north);
\end{pdfdiagram}

\term{有限状态机}{finite-state machine, FSM}由有限状态集合、输入、转换规则和
输出组成。例如一个极简内存读控制器可有：

\[
\text{Idle}\rightarrow\text{Address}\rightarrow\text{Wait}
\rightarrow\text{Complete}\rightarrow\text{Idle}.
\]

软件设备驱动中的状态机和硬件 FSM 使用相同推理方式：先列状态与合法转换，
再确定每个转换的触发条件和超时失败。

\topic{复位决定初始状态，时钟决定推进时刻}
上电时触发器状态未必自动为 0。复位网络把关键状态强制到已知值。复位可以
异步置位、同步释放，也可以完全同步；具体选择取决于时钟是否已经存在、器件
库和跨时钟域要求。

\begin{osgridlongtable}{L{0.20\linewidth}L{0.34\linewidth}L{0.36\linewidth}}
方式 & 优点 & 风险与约束 \\ \hline
\endfirsthead
方式 & 优点 & 风险与约束 \\ \hline
\endhead
异步置位 & 无需时钟即可立即进入复位 & 释放靠近时钟边沿可能破坏 recovery/removal 要求 \\ \hline
同步复位 & 释放天然与时钟对齐 & 时钟停止时不能改变状态，复位要进入数据路径 \\ \hline
异步置位、同步释放 & 兼顾早期强制与安全退出 & 每个时钟域通常需要自己的同步释放链 \\ \hline
\end{osgridlongtable}

CPU 的架构复位状态是大量内部硬件状态机完成复位后的软件可见结果。操作系统
不需要知道每个晶体管怎样复位，但必须严格遵守处理器手册承诺的寄存器、模式
和取指地址。

\topic{从逻辑积木到 ALU}
\term{算术逻辑单元}{arithmetic logic unit, ALU}组合加法、减法、按位逻辑、
移位和比较。操作码经控制逻辑选择一种结果，并产生 zero、carry、sign、
overflow 等条件。

\begin{pdfdiagram}
  \node[os state] (a) {操作数 A};
  \node[os state,below=of a] (b) {操作数 B};
  \node[os block,right=25mm of a,yshift=-10mm] (alu) {ALU\\加减、逻辑、移位};
  \node[os state,right=25mm of alu] (result) {结果};
  \node[os state,below=of result] (flags) {条件标志};
  \node[below=of alu] (op) {操作选择};
  \draw[os bus] (a) -- (alu);
  \draw[os bus] (b) -- (alu);
  \draw[os arrow] (op) -- (alu);
  \draw[os bus] (alu) -- (result);
  \draw[os arrow] (alu) -- (flags);
\end{pdfdiagram}

ALU 本身不会决定“下一条执行什么”。程序计数器、寄存器文件、指令译码、
访存接口和控制状态机共同组成处理器。下一章把这些积木连接成一台能够取指的
完整计算机。

\begin{contractbox}
读完本章，应能从真值表写出简单布尔表达式；手算半加器、全加器和译码窗口；
区分组合逻辑与时序逻辑；解释传播延迟、建立/保持时间、亚稳态和复位；并能
把“CPU 做加法”拆成寄存器提供输入、ALU 形成结果、时钟边沿提交状态三步。
\end{contractbox}
